\documentclass[11pt,a4paper]{article}
\usepackage[utf8]{inputenc}
\usepackage[T1]{fontenc}
\usepackage{lmodern}
\usepackage[portuguese]{babel}
\usepackage[a4paper,margin=2.1cm]{geometry}
\usepackage{xcolor}
\usepackage{graphicx}
\usepackage{titlesec}
\usepackage{tcolorbox}
\usepackage{amsmath}
\usepackage{fancyhdr}
\usepackage{hyperref}
\tcbuselibrary{skins,breakable}

\definecolor{BgCream}{HTML}{FAF8F5}
\definecolor{TextEspresso}{HTML}{4A4238}
\definecolor{NudeTaupe}{HTML}{BFAFA6}
\definecolor{SoftBlush}{HTML}{D9C5B2}

\hypersetup{colorlinks=true, linkcolor=TextEspresso, urlcolor=TextEspresso}

\pagestyle{fancy}
\fancyhf{}
\renewcommand{\headrulewidth}{0pt}
\fancyfoot[C]{\color{NudeTaupe}\small\thepage}
\setlength{\parindent}{0pt}
\setlength{\parskip}{0.32cm}
\setlength{\emergencystretch}{3em}

\newtcolorbox{slidebox}{
  breakable, colback=BgCream, colframe=NudeTaupe,
  boxrule=0.6pt, arc=2pt, left=4pt, right=4pt, top=4pt, bottom=4pt
}

\newcommand{\sebentatitle}[3]{%
  \thispagestyle{empty}
  \begin{center}
  {\color{NudeTaupe}\rule{\textwidth}{0.5pt}}\\[0.5cm]
  {\LARGE\bfseries\color{TextEspresso} Sebenta #1}\\[0.35cm]
  {\large\color{TextEspresso} #2}\\[0.3cm]
  {\normalsize\color{NudeTaupe} Infraestruturas de Computação na Nuvem \quad|\quad #3}\\[0.4cm]
  {\color{NudeTaupe}\rule{\textwidth}{0.5pt}}
  \end{center}
  \vspace{0.4cm}
}

\newcommand{\pagina}[2]{%
  \begin{slidebox}
  \centering
  \includegraphics[width=0.95\textwidth]{#1}
  \end{slidebox}
  \vspace{0.4cm}
  #2
  \clearpage
}

\begin{document}

\sebentatitle{07}{Automatização de Configurações - IaC}{Aula 07}

\section*{Objeto e utilização deste documento}

O presente documento acompanha a apresentação da Aula 07. Cada página reproduz um slide e desenvolve, em texto corrido, a matéria nele condensada. A aula é expositiva e prolonga as anteriores: a abordagem declarativa já conhecida do Kubernetes (Aula 05), a ideia de que a configuração de rede pode ser centralizada e programada (Aula 06) e o isolamento e a imutabilidade das imagens de containers (Aula 04).

A leitura sequencial é a mais proveitosa. A aula parte da tensão entre desenvolvimento e operações que originou o movimento DevOps e das limitações da configuração manual de servidores. Apresenta depois o princípio de Infrastructure as Code e as duas famílias de ferramentas que o concretizam, a gestão de configuração com o Ansible e o provisionamento com o Terraform. Termina com a gestão de segredos, o CI/CD aplicado a infraestrutura e uma síntese de conceitos.

\clearpage

\pagina{slides/slide-01.png}{
A Aula 06 mostrou que a rede de um datacenter pode ser gerida como software, a partir de um ponto central. A pergunta que esta aula coloca é mais geral: o que acontece quando toda a infraestrutura, e não apenas a rede, passa a ser descrita e aplicada como código?

A resposta é a automatização da configuração e do provisionamento, conhecida pela expressão Infrastructure as Code (IaC). A aula apresenta as motivações que a justificam, os princípios que a sustentam e as ferramentas que a tornam prática, em particular o Ansible e o Terraform.
}

\pagina{slides/slide-02.png}{
Os objetivos de aprendizagem organizam o estudo em três eixos. O primeiro é conceptual: justificar a automatização face à gestão manual e distinguir as abordagens declarativa e imperativa, com o conceito de idempotência. O segundo é técnico: descrever o Ansible, explicar o Terraform e o seu state file e diferenciar provisionamento de gestão de configuração.

O terceiro eixo é de prática profissional: aplicar boas práticas de gestão de segredos e de CI/CD à infraestrutura. Ao terminar o estudo deste documento, o estudante deve conseguir explicar cada um destes pontos com as suas próprias palavras e identificar, perante um cenário concreto, qual a ferramenta adequada.
}

\pagina{slides/slide-03.png}{
A primeira metade da agenda trata da motivação e dos fundamentos. Começa pelo DevOps, entendido como resposta organizacional à tensão entre quem desenvolve software e quem opera a infraestrutura, e pelas limitações da configuração manual, que tornam essa tensão dispendiosa.

Seguem-se o princípio de Infrastructure as Code, a distinção entre as abordagens declarativa e imperativa, com a noção de idempotência, e a gestão de configuração com o Ansible.
}

\pagina{slides/slide-04.png}{
A segunda metade da agenda passa do conceito à prática. Apresenta o provisionamento declarativo com o Terraform e o papel do state file, e esclarece a diferença entre provisionamento e gestão de configuração, uma das confusões mais frequentes neste tema.

Termina com dois assuntos transversais a qualquer ferramenta: a gestão de segredos, que impede que credenciais acabem expostas no repositório, e o CI/CD aplicado a infraestrutura, que automatiza a validação e a aplicação das alterações.
}

\pagina{slides/slide-05.png}{
O diagrama resume o percurso da aula como um fluxo único. Uma alteração nasce num repositório Git, no GitHub, sob a forma de commit e de pull request. Um pipeline de CI/CD valida o código, gera um plano e, após revisão, aplica-o. O Terraform provisiona as máquinas e as redes, e o Ansible configura o software que corre nelas, resultando numa infraestrutura convergente e auditável.

Dois elementos acompanham o fluxo. Os segredos são fornecidos ao pipeline apenas em tempo de execução, e não ficam no repositório. A barra inferior representa a reconciliação contínua, em que o Git funciona como fonte de verdade e qualquer divergência do estado real é detetada e corrigida. Cada bloco será desenvolvido nas páginas seguintes.
}

\pagina{slides/slide-06.png}{
Historicamente, as equipas de desenvolvimento (Dev) e de operações (Ops) têm incentivos opostos. As primeiras são avaliadas pela rapidez com que entregam funcionalidades, as segundas pela estabilidade dos sistemas em produção, e cada alteração é, para quem opera, um risco. Esta separação gera passagens de testemunho manuais, ambientes inconsistentes e um lead time elevado entre o commit e a entrada em produção.

A resposta burocrática, feita de aprovações manuais e janelas de alteração rígidas, agrava o problema: aumenta o lead time sem melhorar a qualidade. O DevOps surge como a cultura e o conjunto de práticas que integram as duas funções num único ciclo, apoiado em versionamento, integração e entrega contínuas e automatização.
}

\pagina{slides/slide-07.png}{
O símbolo do infinito representa o ciclo DevOps como um processo contínuo, e não como uma sequência com princípio e fim. O laço do lado do desenvolvimento reúne o planeamento, a escrita do código, a construção (build) e os testes. O laço do lado das operações reúne a disponibilização (release), a implantação (deploy), a operação e a monitorização.

A monitorização alimenta o planeamento seguinte, fechando o ciclo. No cruzamento dos dois laços está a infraestrutura descrita como código, porque é ela que garante que o ambiente onde o software é testado e o ambiente onde é operado são o mesmo, e que ambos podem ser recriados a pedido.
}

\pagina{slides/slide-08.png}{
O DevOps materializa-se em cenários concretos de engenharia de software. Na integração contínua, cada commit desencadeia a compilação e a execução automática de testes. Na entrega e implantação contínuas, as versões seguem para produção de forma frequente e reversível.

As implantações progressivas, como canary e blue-green, e as feature flags limitam o impacto de cada alteração, porque expõem a nova versão a uma fração dos utilizadores ou permitem desligá-la sem nova implantação. A criação de ambientes efémeros por pull request exige infraestrutura provisionada a pedido. A observabilidade, com métricas, logs e alertas, alimenta o planeamento seguinte, e o DevSecOps integra no pipeline a análise de dependências e a deteção de segredos expostos.
}

\pagina{slides/slide-09.png}{
Este slide justifica a presença do tema nesta unidade curricular. Todos os cenários anteriores pressupõem que os ambientes de destino existem, estão atualizados e são equivalentes entre si. Quando a infraestrutura é configurada manualmente, torna-se o elo não automatizado do ciclo e reintroduz as passagens manuais que o DevOps procura eliminar.

A automatização da infraestrutura é, por isso, a condição técnica para que esses cenários possam ser executados de forma contínua, repetível e auditável. É nesse sentido que o DevOps é aqui abordado.
}

\pagina{slides/slide-10.png}{
A configuração manual de servidores tem limitações estruturais. Não escala, porque o esforço cresce linearmente com o número de máquinas. Origina \emph{configuration drift}, a evolução independente de cada servidor, difícil de detetar sem um inventário do estado real. E expõe-se ao erro humano, com passos omitidos ou comandos mal aplicados que geram inconsistências difíceis de diagnosticar.

Há ainda consequências operacionais. Os ambientes de desenvolvimento, teste e produção divergem, pelo que os testes deixam de ser representativos. A recuperação de desastres é lenta, porque recriar a infraestrutura à mão demora horas ou dias, o que se traduz num \emph{recovery time objective} (RTO) elevado. Por fim, a auditoria e a conformidade exigem um registo de quem alterou o quê e quando, que a prática manual raramente fornece.
}

\pagina{slides/slide-11.png}{
O diagrama compara as duas abordagens. À esquerda, um administrador liga-se por SSH a cada servidor e aplica as alterações à mão. Um servidor fica com um passo omitido, outro recebe um \emph{hotfix} manual, e o conjunto diverge ao longo do tempo sem que haja registo do que mudou.

À direita, a definição desejada vive num repositório Git, em ficheiros HCL ou YAML. Um motor de IaC lê essa definição e aplica-a de forma idempotente a todos os servidores, que convergem para o mesmo estado. A diferença não está no objetivo, que é o mesmo, mas no processo: o primeiro depende da memória e da disciplina de quem executa, o segundo é repetível por construção.
}

\pagina{slides/slide-12.png}{
Infrastructure as Code consiste em descrever a infraestrutura, isto é, servidores, redes e serviços, em ficheiros de texto, em vez de a construir através de passos manuais. Esses ficheiros, escritos em linguagens como YAML ou HCL, são versionados em Git, o que dá histórico, revisão por pares e a possibilidade de repor qualquer versão anterior.

Duas propriedades tornam a abordagem útil. A idempotência garante que aplicar a mesma definição várias vezes converge sempre para o mesmo estado. A repetibilidade permite gerar ambientes idênticos em qualquer momento e testar alterações antes de chegarem a produção, com verificações de sintaxe (\emph{lint}) e execuções simuladas (\emph{dry-run}). O código funciona ainda como documentação atualizada e reduz o tempo de provisionamento de dias para minutos.
}

\pagina{slides/slide-13.png}{
O IaC é um princípio, e não uma ferramenta. Estabelece que a infraestrutura deve ser descrita, versionada e aplicada como código, mas não diz como. A sua concretização envolve dois problemas distintos, e cada um tem as suas ferramentas.

O primeiro é garantir que os recursos existem, como máquinas virtuais, redes e armazenamento. É o provisionamento, de que o Terraform é o exemplo estudado. O segundo é definir o que corre dentro de cada recurso, nomeadamente pacotes, ficheiros de configuração e serviços. É a gestão de configuração, de que o Ansible é o exemplo estudado. Em ambos se aplicam os mesmos fundamentos, a descrição declarativa do estado e a idempotência.
}

\pagina{slides/slide-14.png}{
O Ansible é uma ferramenta open source de gestão de configuração e automatização. Aplica, por SSH e sem necessitar de um agente instalado, o estado descrito em ficheiros YAML denominados playbooks a máquinas que já existem.

O Terraform é uma ferramenta open source de provisionamento declarativo. A partir de código HCL, cria, altera e destrói recursos de infraestrutura, como máquinas virtuais, redes e armazenamento, recorrendo às APIs dos \emph{providers}, que são os plugins específicos de cada plataforma. As duas ferramentas respondem, portanto, a perguntas diferentes: o Terraform decide que recursos existem, o Ansible decide o que neles corre.
}

\pagina{slides/slide-15.png}{
Na abordagem imperativa descreve-se a sequência de passos a executar, como acontece num script shell que instala um pacote, copia um ficheiro e reinicia um serviço. O resultado depende do estado inicial da máquina, e repetir o script pode produzir efeitos indesejados, como acrescentar duas vezes a mesma linha a um ficheiro.

Na abordagem declarativa descreve-se o estado final pretendido, e é a ferramenta que calcula as ações necessárias para lá chegar. É o mesmo princípio do Kubernetes, estudado na Aula 05. Esta abordagem torna possível a idempotência, segundo a qual reaplicar a definição converge sempre para o mesmo estado. As ferramentas distinguem o que mudou (\emph{changed}) do que já estava conforme (\emph{ok}), o que torna cada execução auditável.
}

\pagina{slides/slide-16.png}{
O Ansible comunica por SSH, na porta 22, e não exige agente nos hosts geridos. Para cada task, o módulo correspondente é transferido para o host, executado com Python e removido no final. O requisito mínimo no destino é, por isso, um servidor SSH e um interpretador Python.

Os playbooks são escritos em YAML e descrevem o estado desejado de forma declarativa. O \emph{inventory} lista os hosts geridos, agrupados por função, por exemplo num grupo de servidores web. Os módulos, como \texttt{apt}, \texttt{copy} e \texttt{service}, são idempotentes, pelo que um playbook pode ser executado repetidamente. O Ansible é mantido pela Red Hat e distribuído como software open source.
}

\pagina{slides/slide-17.png}{
Um projeto Ansible organiza-se em torno de alguns componentes. Um playbook é um ficheiro YAML que associa um conjunto de tasks a um grupo de hosts. Cada task é uma ação individual, executada por um módulo.

Um handler é uma task especial, executada apenas quando outra notifica que algo mudou, por exemplo para reiniciar um serviço depois de alterar o seu ficheiro de configuração. Uma role agrupa tasks, variáveis, templates e ficheiros numa unidade reutilizável. Um template Jinja2 é um ficheiro de configuração parametrizado por variáveis, instanciado para cada host. Os \emph{facts} são informação recolhida automaticamente sobre cada host, como o sistema operativo, a memória e os endereços IP, e podem ser usados em condições e templates.
}

\pagina{slides/slide-18.png}{
O diagrama mostra como estes componentes se relacionam. No control node, o comando \texttt{ansible-playbook}\allowbreak{} lê o inventory, que define os hosts alvo, e o playbook, que define o que aplicar. O playbook usa tasks, que invocam módulos, e as tasks podem ser organizadas em roles, que incluem templates Jinja2.

Quando uma task altera algo, notifica o handler associado, que é executado no fim, e apenas uma vez. As tasks são enviadas por SSH para os hosts geridos, onde são executadas e onde são recolhidos os facts. O control node mantém todo o estado desejado, e os hosts apenas executam.
}

\pagina{slides/slide-19.png}{
O exemplo mostra um playbook mínimo para servidores web. A primeira task usa o módulo \texttt{apt} para garantir que o pacote Nginx está instalado (\texttt{state: present}). A segunda copia o ficheiro de configuração a partir de um template Jinja2 e, se o conteúdo mudar, notifica o handler.

O handler reinicia o serviço, mas só quando notificado. Numa segunda execução sem alterações, todas as tasks reportam \emph{ok} e o handler não é executado. Este comportamento ilustra a idempotência em prática: o playbook descreve o estado, e não os passos, e só atua quando há diferença.
}

\pagina{slides/slide-20.png}{
O Terraform é uma ferramenta declarativa para provisionar recursos de infraestrutura. Os \emph{providers} são plugins que traduzem o código em chamadas às APIs de clouds ou de plataformas específicas, como o Proxmox. O state file regista o que foi efetivamente criado, para que possa ser comparado com o estado desejado.

O fluxo de trabalho assenta nos comandos \texttt{plan} e \texttt{apply}. O primeiro mostra as alterações previstas e o segundo executa-as, o que permite rever o impacto antes de o aplicar. A linguagem usada é a HCL, \emph{HashiCorp Configuration Language}, desenhada para descrever recursos de forma legível.
}

\pagina{slides/slide-21.png}{
O exemplo declara dois blocos principais. O bloco \texttt{terraform} indica o provider necessário, neste caso o do Proxmox. O bloco \texttt{resource} descreve o recurso a criar, uma máquina virtual com dois núcleos e 2048~MB de memória, cujo nome é calculado a partir de um índice.

A variável \texttt{vm\_count} define quantas máquinas existem. Alterar o seu valor de 2 para 3 e executar \texttt{terraform plan} mostra a criação de uma única máquina adicional, porque o Terraform compara o estado desejado com o que já existe e calcula apenas a diferença. Não é necessário descrever como criar a máquina nem escrever lógica para evitar duplicados.
}

\pagina{slides/slide-22.png}{
O diagrama mostra a arquitetura interna do Terraform. O código \texttt{.tf} é lido pelo núcleo (\emph{core}), que constrói um grafo de dependências entre recursos. O núcleo consulta e atualiza o state file, que reside num backend remoto com bloqueio, e determina o plano de ações.

Esse plano é entregue ao provider, um plugin que comunica com o núcleo por gRPC, e que por sua vez chama a API REST da plataforma, por HTTPS, para criar, alterar ou destruir os recursos reais. O state file liga, assim, os recursos descritos no código aos recursos reais, e por isso deve ser tratado como um dado sensível.
}

\pagina{slides/slide-23.png}{
O state file regista o mapeamento entre os recursos descritos no código e os recursos reais. É este mapeamento que permite ao Terraform determinar, em cada execução, o que já existe, o que falta criar e o que deve ser destruído.

Em trabalho de equipa, o estado deve residir num backend remoto partilhado, com \emph{locking} para impedir execuções simultâneas que o corrompam. Como pode conter dados sensíveis em texto simples, deve ser tratado como um segredo. Alterações manuais feitas fora do Terraform introduzem drift e serão revertidas na aplicação seguinte, o que reforça a regra de que toda a alteração deve passar pelo código.
}

\pagina{slides/slide-24.png}{
O diagrama mostra o ciclo \emph{plan} e \emph{apply}. O comando \texttt{plan} lê o código, que descreve o estado desejado, e o state file, que descreve o que já foi criado, compara-os e mostra as alterações previstas. Nada é modificado nesta fase.

Se o operador confirmar, o \texttt{apply} executa as alterações na infraestrutura real, neste caso o Proxmox, e atualiza o state file. A separação em duas fases é uma salvaguarda importante. Permite rever, e num pipeline de CI/CD aprovar, o que vai acontecer antes de acontecer.
}

\pagina{slides/slide-25.png}{
Este slide fixa a distinção entre as duas camadas. O provisionamento decide que máquinas e recursos existem, e o Terraform é o exemplo. A gestão de configuração decide o que corre dentro de cada máquina, e o Ansible, tal como o Puppet, são exemplos.

O fluxo habitual combina as duas: o Terraform cria a máquina virtual e o Ansible instala e configura o software. As camadas são complementares e não concorrentes, e confundi-las é a origem mais comum de dúvidas sobre qual ferramenta usar para uma determinada tarefa.
}

\pagina{slides/slide-26.png}{
Esta página mostra como as tecnologias estudadas ao longo da unidade curricular se encadeiam. O Terraform provisiona as máquinas, a rede e o armazenamento, incluindo os nós de um cluster. O Ansible configura esses nós e instala o container runtime. O Docker empacota a aplicação numa imagem imutável, publicada num registry pelo pipeline de CI/CD. O Kubernetes executa essa imagem de forma declarativa, a partir de manifests YAML.

A ordem apresentada é uma simplificação, porque a construção da imagem não depende de o Ansible ter corrido. O essencial é que cada camada é descrita em código, versionada no Git e aplicada pelo mesmo pipeline, o que garante consistência entre a infraestrutura de base e a aplicação.
}

\pagina{slides/slide-27.png}{
O código de infraestrutura é versionado em Git, mas passwords, chaves e tokens não podem ser guardados no repositório em texto simples. Este é o principal desafio de segurança do IaC.

O Ansible Vault cifra ficheiros de variáveis, o que permite guardá-los no repositório. O HashiCorp Vault e serviços equivalentes armazenam os segredos de forma centralizada e fornecem-nos às ferramentas apenas em tempo de execução. As variáveis de ambiente podem ser injetadas pelo sistema de CI/CD, sem nunca serem escritas no repositório. Um segredo que tenha sido submetido (\emph{commit}) deve ser considerado comprometido e rodado, porque permanece no histórico do Git mesmo depois de apagado do ficheiro.
}

\pagina{slides/slide-28.png}{
Aplicam-se ao código de infraestrutura os mesmos princípios de integração e entrega contínuas usados no software. O pipeline valida a sintaxe, gera o plano, submete-o a revisão e só então executa o \emph{apply}.

O diagrama mostra estas etapas em sequência: o \emph{push} ou pull request inicia o pipeline, seguem-se as validações de formato e sintaxe, a geração do plano como artefacto, a revisão por pares e, por fim, a aplicação num runner de CI. Se qualquer etapa falhar, o pipeline termina com um código de saída diferente de zero e a alteração não avança. Os testes automatizados antecedem qualquer alteração em produção, o que reduz o risco de alterações manuais não controladas.
}

\pagina{slides/slide-29.png}{
A tabela reúne os conceitos técnicos da aula. A idempotência significa que reaplicar a definição não altera o estado. O \emph{configuration drift} é a divergência entre o estado real e o declarado, por exemplo uma edição manual de um ficheiro de configuração. Na abordagem declarativa especifica-se o estado final e o motor calcula a diferença.

O state file é o mapeamento entre recursos declarados e reais. A distinção entre agentless e agent-based opõe o modelo de \emph{push} por SSH, do Ansible, ao de \emph{pull} por um agente residente, do Puppet e do Chef. A imutabilidade substitui a instância em vez de a alterar no local, como acontece com imagens Packer ou Docker. A reconciliação é o ciclo contínuo de observar, comparar e corrigir, que ferramentas como o Argo CD e o Flux aplicam.
}

\pagina{slides/slide-30.png}{
A primeira parte da síntese recolhe os fundamentos. O DevOps responde à tensão entre desenvolvimento e operações através de cultura, automatização e versionamento. A configuração manual conduz a drift, erros e falta de reprodutibilidade.

O IaC trata a infraestrutura como código versionado, idempotente e repetível. O Ansible, agentless e baseado em SSH e YAML, contrasta com o Puppet e o Chef, que mantêm um agente em cada host e verificam periodicamente o estado, uma opção adequada a parques de grande dimensão.
}

\pagina{slides/slide-31.png}{
A segunda parte da síntese recolhe as ferramentas e as práticas. O Terraform provisiona recursos de forma declarativa, com providers, state e o ciclo plan e apply, e o state file é crítico e sensível. O provisionamento e a gestão de configuração são camadas complementares.

Os segredos nunca vão para o repositório em texto simples. O CI/CD automatiza o teste e a aplicação das alterações, tendo o Git como fonte de verdade. A Aula 08 retoma esta base: a automatização que aqui permite criar e configurar máquinas a pedido é a condição para o \emph{auto-scaling}.
}

\pagina{slides/slide-32.png}{
Para aprofundar a matéria, recomendam-se as obras de Morris sobre Infrastructure as Code, de Brikman sobre o Terraform e de Hochstein e Moser sobre o Ansible, todas da O'Reilly. A documentação oficial das duas ferramentas é a referência atualizada para a sintaxe e para os módulos e providers disponíveis.
}

\pagina{slides/slide-33.png}{
Este é o espaço para questões e comentários. Sugere-se que o estudante, antes da Aula 08, consiga explicar com palavras suas três coisas: a diferença entre provisionamento e gestão de configuração, o que significa dizer que uma operação é idempotente e porque é que o state file do Terraform deve ser tratado como um segredo.
}

\end{document}
